% !PRINTED 547
\weberpage{547}
%<<CONT
equation whose roots are the class invariants belonging to the second genus, and which are therefore conjugate imaginary, while the roots of (5) are real. From this there follows the determination of quantity:
\[
(a+b\sqrt p)^2 > 4 > (a-b\sqrt p)^2;
\]
and so $a$ and $b$ must have the same sign. But since $x$ is positive by (4), $a$ and $b$ must both be positive.

In order actually to find $a$ and $b$, one needs only to compute the expression on the left-hand side of (5), by (4), to a few decimals, and for this it is amply sufficient to put
\[
f(\sqrt{-m}) = f_1(\sqrt{-m}) = e^{\frac{\pi\sqrt m}{24}}
\]
and to compare the decimals so found with the decimals of $\sqrt p$, in order to obtain $b$ at once and thereupon $a$. The calculation is exceedingly simple and gives the following results:
\[
\begin{aligned}
m &= 14, & x + \frac1x &= 1 + \sqrt2,\\
m &= 17, & x + \frac1x &= \frac{1+\sqrt{17}}{2},\\
m &= 34, & x + \frac1x &= \frac{3+\sqrt{17}}{2},\\
m &= 46, & x + \frac1x &= 3 + \sqrt2,\\
m &= 49, & x + \frac1x &= 2 + \sqrt7,\\
m &= 73, & x + \frac1x &= \frac{5+\sqrt{73}}{2},\\
m &= 82, & x + \frac1x &= \frac{15+\sqrt{41}}{2},\\
m &= 97, & x + \frac1x &= \frac{9+\sqrt{97}}{2},\\
m &= 142, & x + \frac1x &= 9 + 5\sqrt2,\\
m &= 193, & x + \frac1x &= 13 + \sqrt{193}.
\end{aligned}
\]

\begin{center}\rule{3cm}{0.4pt}\end{center}
